\FloatBarrier
\Needspace{11\baselineskip}
\section{Discussion}
\label{sec:discussion}

\subsection{What a clean sanitizer run means for a program with a CMA}

A clean sanitizer run provides limited evidence about the safety of objects inside an unannotated
pool. In our corpus, the tools detected errors only when an access crossed the pool's boundary or
occurred after the backing pool had been freed. Fuzzing under ASan may therefore execute invalid
accesses within a pool without reporting them.
Three practical consequences follow.
\begin{itemize}
\item \emph{Annotate, and choose the tool by the bug class.} Adding annotation macros at 14 call
  sites in Arena and 24 in Free list recovers most bugs in our corpus. ASan's annotations increased
  median runtime by approximately 5--10\% on our allocator microbenchmark, but poisoning cannot
  detect invalid or double frees, overflows into live neighbours, or stale-pointer accesses after
  storage reuse. Memcheck's memory-pool interface also detects invalid and double frees, with
  annotated median runtimes 1.8--2.2$\times$ those of the unannotated builds. HWASan's tags detected
  neighbour overflows and stale-pointer accesses after reuse in our tested cases, but its annotation
  interface has 16-byte precision and does not detect invalid or double frees. The tools therefore
  provide complementary coverage (\S\ref{sec:q1}, \S\ref{sec:perf}).
\item \emph{Annotation needs the allocator's cooperation.} To re-poison an object on free, an
  allocator must know which bytes belong to it. Our Arena uses a side table to track object sizes
  and lifetimes; a bump allocator without equivalent tracking cannot reliably mark individually
  freed objects as inaccessible. Free list must temporarily make the link words inside freed slots
  accessible while reading or updating them. Annotation therefore requires changes to allocator
  bookkeeping as well as calls to the tool's interface.
\item \emph{A configurable CMA can be disabled in a test build.} In our SQLite rewrite configuration,
  disabling memsys5, lookaside and the page-cache bulk pool exposes the tested objects as individual
  \code{malloc} allocations. This requires no changes to SQLite's allocator implementation and
  restores \code{malloc}-level detection, but it can change performance and allocation behaviour.
  In our corpus, bugs that require immediate reuse or bulk deallocation become unreproducible when
  those behaviours are removed (\S\ref{sec:q1-rewrite}).
\end{itemize}

\subsection{Do the results transfer to real bugs?}
\label{sec:tier2}

To see whether the corpus's verdicts transfer to real defects, we ported four bugs from Capstone's
SQLite corpus~\cite{capstone} that are live in SQLite 3.22.0: two uses after free (a connection read
after \code{sqlite3\_close}, a stale R-tree node) and two fuzzer-found over-reads of page buffers. The
drivers are Capstone's, unchanged, and SQLite itself is not simplified, since cutting the engine
would move the objects. Each bug runs in three configurations. In the first, memsys5 serves every allocation from a 256~KiB
static arena. The second is SQLite as shipped, where the connection and the R-tree node are plain
\code{malloc} blocks and pages come from the page cache's bulk block. The third is the rewrite (no
memsys5, lookaside and bulk block off), which makes every object a \code{malloc} block and is each
bug's positive control. A real bug has no negative twin, so the tools' positive controls and the rewrite
gate the run. Table~\ref{tab:real} gives the formal run on the lab server.

\begin{table}[!htbp]
\centering\small
\setlength{\tabcolsep}{4.5pt}
\begin{tabular}{rl*{6}{c}}
\toprule
 & & \multicolumn{3}{c}{ASan} & \multicolumn{3}{c}{Valgrind} \\
\cmidrule(lr){3-5}\cmidrule(lr){6-8}
\# & bug & memsys5 & shipped & rewrite & memsys5 & shipped & rewrite \\
\midrule
03 & use after free: connection read after close & \cM & \cD & \cD & \cM & \cD & \cD \\
19 & use after free: stale R-tree node & \cMs & \cD & \cD & \cD$^{\dagger}$ & \cD & \cD \\
29 & over-read of a page (\code{sqlite3Get4byte}) & \cM & \cD$^{\ddagger}$ & \cD & \cM & \cD$^{\ddagger}$ & \cD \\
31 & over-read of a page (\code{statDecodePage}) & \cD$^{\ddagger}$ & \cD$^{\ddagger}$ & \cD & \cM & \cD$^{\ddagger}$ & \cD \\
\bottomrule
\end{tabular}

\caption{Four real SQLite 3.22.0 bugs on the lab server (\code{results/2026-11-16-lab-real}). Shipped
= SQLite's default configuration. $\dagger$: Memcheck's report is an invalid read at address
\code{0x2} just before the process dies of \code{SIGSEGV}, the same crash that ASan records as
M$^{*}$, and it does not concern the stale node. $\ddagger$: the over-read left its container (the arena
or the page cache's bulk block) and is reported against a redzone, a freed block or unallocated
heap, never against the page.}
\label{tab:real}
\end{table}

\paragraph{The temporal results transfer directly.} memsys5 hides both uses after free from both
tools, while in the shipped configuration and the rewrite, where the freed objects are
\code{malloc} blocks, both tools report the read and name the freed block. In the R-tree case the
stale read itself is silent. A pointer loaded from the stale node is \code{0x2}, and the program
then faults on the zero page. ASan's signal handler reports the fault (missed, crashed), and Memcheck
reports an invalid read at \code{0x2} that the harness counts as detected. Neither report concerns
the node, and we read both as misses.

\paragraph{The spatial results depend on the stride, as in cases 3 and 4.} Both over-reads go far:
under the rewrite, ASan reports case 29's first bad read 253 bytes past a 784-byte page block,
and Memcheck reports case 31's 51~KB into unallocated heap. Inside memsys5, case 29 stays in
memory the tools consider valid and is missed by both. Case 31 is missed by Memcheck, which has no
redzones around global data, and reported by ASan only because it lands 34{,}968 bytes past the end
of the arena, in the arena's own redzone. In the shipped configuration both reads leave the bulk
block, and both tools report them against heap redzones or other blocks. Every spatial detection
outside the rewrite is therefore a detection of the container, not of the page, and a CMA hides a
real over-read whenever the read stays inside it. A development run on macOS had detected case 29
inside memsys5, 1.9~KB past the arena: like case 4 of the corpus, the verdict of a far stride
depends on the address-space layout.

\subsection{Limitations}
\label{sec:limits}

\begin{itemize}
\item \textbf{UBSan and LowFat are not in the matrix.} At \code{-O0} clang's \code{object-size}
  check has no effect (the compiler warns so), and \code{bounds} needs a statically known array
  length, which a pointer returned by an allocator does not have; neither positive control is
  detected at \code{-O0}, and at \code{-O2} only the overflow is. These UBSan checks do not provide
  general per-object bounds and lifetime tracking for CMA objects; no full UBSan matrix was produced.
  Our working LowFat configuration used its LLVM 4.0 compiler with glibc 2.27 in an Ubuntu 18.04
  container. There it missed both positive controls (a 64-byte block's one-byte overflow falls inside its size
  class, and it has no temporal checks) while catching an overflow that crosses the size-class bound.
  Both tools would need different positive controls to enter the harness.
\item \textbf{Synthetic corpus.} Sixteen single-threaded programs isolate mechanisms; they do not
  estimate how often each mechanism occurs in real code. The real-bug check covers four SQLite defects and tests whether the mechanisms transfer, not how common they are.
\item \textbf{One configuration per tool.} One recorded option set, \code{-O0} for detection, one run per cell.
  HWASan is probabilistic, and one cell was affected; case 4 under ASan depends on the address-space
  layout.
\item \textbf{Different machines.} HWASan ran on another architecture, and the performance numbers
  come from two machines different from the detection server. We never combine numbers across
  machines. The three SQLite versions (3.53.3 in the corpus, 3.50.4 in \code{speedtest1}, 3.22.0 in
  the real bugs) are reported separately.
\item \textbf{Experimental scope.} Detection used native Linux environments, with LowFat feasibility
  checks in Docker. The study covers 16 controlled bug programs, four real SQLite
  cases, and performance measurements on SQLite, Lua and an allocator microbenchmark. PostgreSQL
  and CPython were not evaluated.
\end{itemize}
